\section{All-order factorization defects: local residues and convergent terms}
\label{hqj:sec:local}

Define initially in their absolute-convergence half-planes
\begin{align}
Z_W(s;a,b)&=\sum_{n\ge0}W(x_n)(x_n^2+b^2)^{-s},\label{hqj:eq:ZW}\\
F_{W,\pm}(s;a,b)&=\sum_{n\ge0}W(x_n)(x_n\pm ib)^{-s}.
\end{align}
Their initial conditions are respectively $\Rea s>(D+1)/2$ and
$\Rea s>D+1$. Introduce the common-scale difference
\begin{align}
\mathcal D_W(s)&=\tfrac12\bigl(F_{W,+}(2s)+F_{W,-}(2s)\bigr)-Z_W(s),
\label{hqj:eq:defect-germ}\\
\Delta_r(W;a,b)&=\mathcal D_W^{(r)}(0)
=2^{r-1}\bigl(F_{W,+}^{(r)}(0)+F_{W,-}^{(r)}(0)\bigr)-Z_W^{(r)}(0).
\label{hqj:eq:defect}
\end{align}
The factors of two are part of the definition and are essential from
second order onward.

Set
\begin{equation}\label{hqj:eq:logs}
\ell_\pm(x)=\log(1\pm ib/x),\qquad
L(x)=\log(1+b^2/x^2),\qquad
D_b(x)=\ell_+(x)-\ell_-(x)=2i\arctan(b/x).
\end{equation}
The letter $D_b$ is a logarithmic difference, not the degree of $W$.
Pointwise, the summand of \eqref{hqj:eq:defect-germ}, after removing
$W(x)x^{-2s}$, is
\begin{equation}\label{hqj:eq:cosh}
e^{-sL(x)}\bigl(\cosh(sD_b(x))-1\bigr)
=\sum_{k\ge2}h_k(x)s^k,
\end{equation}
where
\begin{equation}\label{hqj:eq:hk}
h_k(x)=\sum_{q=1}^{\lfloor k/2\rfloor}
\frac{D_b(x)^{2q}(-L(x))^{k-2q}}{(2q)!(k-2q)!}.
\end{equation}
In particular,
\[
h_2=\tfrac12D_b^2,\qquad h_3=-\tfrac12LD_b^2,
\qquad h_4=\tfrac14L^2D_b^2+\tfrac1{24}D_b^4.
\]
Only even powers of $x^{-1}$ occur in every $h_k$.

\subsection{An explicit finite-subtraction operator}
For fixed $k$, let
\begin{equation}\label{hqj:eq:Sk}
S_k(x)=\sum_{e\ge-1}a_{k,e}x^e
\end{equation}
be the finite part of the expansion of $W(x)h_k(x)$ at infinity containing
all exponents at least $-1$. Put $R_k=a_{k,-1}$, with missing coefficients
understood to be zero. Then $W(x)h_k(x)-S_k(x)=O(x^{-2})$.
Define
\begin{align}
\mathcal H_k(t)={}&\sum_{n\ge0}\bigl(W(x_n)h_k(x_n)-S_k(x_n)\bigr)x_n^{-t}
+\sum_{e\ge-1}a_{k,e}\zeta(t-e,a),\label{hqj:eq:Hk}\\
\mathcal H_{k,\mathrm{reg}}(t)={}&\mathcal H_k(t)-R_k/t.
\end{align}
The series is normally convergent for $\Rea t>-1$. In a neighborhood of
zero the second line is holomorphic. This construction fixes the
regularizing coordinate to be $x_n=n+a$.

\begin{theorem}[Complete local/nonlocal decomposition]\label{hqj:thm:local}
The functions in \eqref{hqj:eq:ZW} and the two linear-factor functions are
holomorphic at $s=0$ after meromorphic continuation. For every $r\ge1$,
\begin{equation}\label{hqj:eq:all-local}
\boxed{\quad
\frac{\Delta_r(W;a,b)}{r!}
=\frac{R_{r+1}}2+
\sum_{k=2}^r\frac{2^{r-k}}{(r-k)!}
\mathcal H_{k,\mathrm{reg}}^{(r-k)}(0).
\quad}
\end{equation}
Every regular derivative in this formula has the ordinary convergent
representation
\begin{align}
\mathcal H_{k,\mathrm{reg}}^{(j)}(0)
={}&(-1)^j\sum_{n\ge0}\bigl(W(x_n)h_k(x_n)-S_k(x_n)\bigr)\log^j x_n
\notag\\
&+\sum_{e\ne-1}a_{k,e}\zeta^{(j)}(-e,a)
+a_{k,-1}(-1)^j\gamma_j(a).
\label{hqj:eq:regular-derivative}
\end{align}
The coefficients $a_{k,e}$ are finite polynomials in $b$ and the
coefficients of $W$. No unperformed divergent sum is present on the right.
\end{theorem}
\begin{proof}
Choose a finite head so that the remaining $x_n$ exceed $2b$. On the tail,
the inverse-$x$ expansions of the logarithms and exponentials converge
uniformly on compact sets of $s$ near zero. Subtract their terms through
inverse degree $D+1$. The resulting summand is locally uniformly
$O(x_n^{-2+2|\Rea s|})$, up to a fixed polynomial in $\log x_n$ after
any fixed number of derivatives. It is therefore normally summable for
$|\Rea s|<1/4$. The subtracted terms are finite multiples of
$\zeta(2s+j-d,a+N)$.

The only possible pole at zero is the term with $j-d=1$.
For each individual spectrum its inverse-power coefficient vanishes at
least linearly in $s$, canceling that simple pole. The degree-zero
coefficient only gives Hurwitz values at nonpositive integers and is
regular. The finite head is entire in $s$. This proves holomorphy of
the individual spectra at zero.

For the difference, \eqref{hqj:eq:cosh} starts at $s^2$. The same
subtraction argument allows coefficient extraction from
$\sum_{k\ge2}s^k\mathcal H_k(2s)$. Since
\[
\mathcal H_k(2s)=\frac{R_k}{2s}
+\sum_{j\ge0}\frac{2^j\mathcal H_{k,\mathrm{reg}}^{(j)}(0)}{j!}s^j,
\]
the coefficient of $s^r$ is exactly \eqref{hqj:eq:all-local}.
Only $k\le r+1$ can contribute. The finite head and the tail can be
recombined in \eqref{hqj:eq:Hk}; its ordinary remainder converges even
when $b>a$. Finally differentiating that normally convergent remainder
and using \eqref{hqj:eq:stieltjes} gives
\eqref{hqj:eq:regular-derivative}.
\end{proof}

\subsection{The first two orders, and what parity really proves}
At the first order, Theorem~\ref{hqj:thm:local} recovers the predecessor's
local formula:
\begin{align}
\Delta_1(W;a,b)
&=\tfrac14[x^{-1}]W(x)D_b(x)^2\notag\\
&=\sum_{j\ge1}w_{2j-1}\frac{(-1)^j b^{2j}}j
\left(H_{2j}-\tfrac12H_j\right).
\label{hqj:eq:first-defect}
\end{align}
The sum is finite and independent of $a$. Its vanishing for every $b>0$
is equivalent to $W$ being even. For example
\[
\Delta_1(1)=0,\qquad \Delta_1(x)=-b^2,\qquad
\Delta_1(x^3)=\tfrac23b^4.
\]
These are cited predecessor results, also recovered here.

For the second derivative we obtain instead
\begin{equation}\label{hqj:eq:second-defect-reg}
\Delta_2(W)=\Reg_x\sum_{n\ge0}W(x_n)D_b(x_n)^2
-\tfrac12[x^{-1}]W(x)L(x)D_b(x)^2,
\end{equation}
where $\Reg_x$ is precisely \eqref{hqj:eq:Hk} with the constant Laurent
term extracted at $t=0$. Evenness of $W$ removes every $R_k$, because
an even inverse-power expansion cannot contain $x^{-1}$. It does
\emph{not} remove the regular parts in \eqref{hqj:eq:all-local}.

\begin{corollary}[An ordinary convergent second defect]\label{hqj:cor:atan}
For every $a,b>0$,
\begin{equation}\label{hqj:eq:atan-square}
\boxed{\Delta_2(1;a,b)=-4\sum_{n=0}^{\infty}
\arctan^2\!\left(\frac{b}{n+a}\right).}
\end{equation}
In particular this defect is strictly negative.
\end{corollary}
\begin{proof}
For $W=1$, $D_b^2=O(x^{-2})$, so no subtraction or residue occurs.
Insert $D_b=2i\arctan(b/x)$ in \eqref{hqj:eq:second-defect-reg}.
\end{proof}
The sign is used to give an exact counterexample, not as the principal
research target. The predecessor correctly warned against higher-jet
factorization; \eqref{hqj:eq:atan-square} makes its missing term explicit.

\begin{theorem}[Universal higher-order rigidity]\label{hqj:thm:rigidity}
Fix $a>0$ and an integer $r\ge2$. For $W\in\C[x]$,
\[
\Delta_r(W;a,b)=0\quad\hbox{for every }b>0
\quad\Longleftrightarrow\quad W=0.
\]
Thus even weights are the exact first-order kernel, but the kernel at
each fixed higher order is trivial when vanishing is required for all
positive masses.
\end{theorem}
\begin{proof}
Suppose $W\ne0$. There is an $n\ge0$ with $W(x_n)\ne0$. Continue the
finite-head/tail expression in the complex $b$-plane along a path from
the positive real axis to a punctured neighborhood of $b_0=ix_n$,
avoiding the other points $\pm ix_k$. Choose the head large enough that
the remaining tail has $x_k>|b_0|+1$. Its uniformly convergent
inverse-power representation is holomorphic near $b_0$. All other
head terms are also holomorphic there on the continued local branches.

For the $n$th term, $x_n+ib$ has a simple zero at $b_0$, while
$x_n-ib$ is nonzero. Formula \eqref{hqj:eq:defect} gives the coefficient
of the highest logarithmic power $\log^r(b-b_0)$ as
\[
(-1)^r\bigl(2^{r-1}-1\bigr)W(x_n),
\]
which is nonzero. Lower powers of that logarithm and holomorphic terms
cannot cancel it. If the original function vanished on the entire
positive real axis, its analytic continuation would vanish along the
chosen path by the identity theorem, contradicting this local
singularity. The converse for $W=0$ is immediate.
\end{proof}
This is an analytic vanishing theorem for a specified family. It is not
a theorem about algebraic independence of numerical special values.
